% !TeX root = ../教材题库_全部条目_转写工作稿.tex
% Source: XJ-B1-C04-TH-P119-Y0670
% 第4章 / 4.3 对数函数 / 教材 p.119 / PDF 第126页
% OCR未校对：公式与文字均可能有误；未完成数学TeX转写；禁止直接发给学生
% 原题及相邻原解截图（核对用）：../题目原图/XJ-B1-C04-TH-P119-Y0670-01.webp
% 本文件只在转写完成后作为题干使用；OCR 条目尚不是可靠题干。
\begin{bankitem}{p.119 思考（含紧邻说明）}{第4章 / 4.3 对数函数 / 教材 p.119 / PDF 第126页}{XJ-B1-C04-TH-P119-Y0670}{OCR未校对：公式与文字均可能有误；未完成数学TeX转写；禁止直接发给学生}
\begin{ocrdraft}
[教材p.119 / PDF第126页]
\par
？思考如何将logbN(b>0且b1，N>0）转换成以鼠a>0且磊乒D为底的对
\par
数呢？
\par
由N=D04N(N>O，乃>0且bI)可得，当a>0且1时
\par
logaN=Iog\ocrglyph{201E}N）=logbN
\par
logab，
\par
所以
\par
IogaN
\par
（N>0，a>O且a1，乃>0且b\#l).
\par
这个公式叫作对数的换底公式．
\par
一般地，对于底数不是10或的对数，利用换底公式，我们可以将其转换成以
\par
lnN
\par
10或为底的对数，即IWN一或IWN=
\par
lg磊
\par
ln磊
\end{ocrdraft}
\end{bankitem}
